% ----------------------------------------------------------
% Fundamentação Teórica
% ----------------------------------------------------------
\chapter{Fundamentação Teórica} \label{cha:fundamentacao}

\section{JOGOS DIGITAIS}

\subsection{Indústria de Jogos Digitais}

A indústria de jogos digitais se tornou como um dos segmentos mais dinâmicos e rentáveis do mercado global de entretenimento. Estimativas apontam que o público global de jogadores alcançou aproximadamente 3,6 bilhões de pessoas em 2025, movimentando receitas anuais de US\$ 188,8 bilhões. Desse montante, os dispositivos móveis continuam liderando o volume financeiro global (55\%), seguidos pelos consoles (24\%) e computadores (21\%) \cite{newzoo2025}. 

\textcolor{red}{*Crie um gráfico de pizza que apresente esses dados. É bom que quebra um pouco o volume de texto do tcc.}

Essa expansão foi impulsionada pela popularização massiva dos \textit{smartphones}, pelo avanço da conectividade em rede e por mudanças comportamentais aceleradas no período pós-pandêmico, consolidação que transformou os jogos de um entretenimento de nicho em um ecossistema econômico e cultural altamente relevante \cite{mbakirtzis2023}. 

No cenário nacional, o mercado brasileiro tem demonstrado grande preparo e vocação exportadora. Mapeamentos do setor indicam a existência de mais de mil estúdios de desenvolvimento ativos espalhados pelo país. Uma das principais características do ecossistema nacional é sua forte inserção global onde cerca de 58\% das desenvolvedoras brasileiras realizam vendas internacionais, impulsionadas pela distribuição digital em plataformas globais (como \textit{Steam}, \textit{Google Play} e \textit{App Store}) e pela alta qualidade na prestação de serviços de co-desenvolvimento \cite{abragames2023}. 

Paralelamente ao crescimento financeiro, a dinâmica de consumo dos jogos passou por uma profunda transformação estrutural. A excelência gráfica isolada deu lugar ao conceito de "jogos como plataforma", nos quais o sucesso comercial de um título depende primordialmente da formação de comunidades ativas, do engajamento social e do incentivo ao conteúdo gerado pelos próprios usuários (\textit{User-Generated Content} - UGC). Como o público contemporâneo consome e descobre novos títulos por meio de criadores de conteúdo e redes sociais, o ciclo de vida dos jogos tornou-se contínuo e altamente dependente do escuta ativa dos jogadores \cite{bain2025}. 


\subsection{Desenvolvimento Colaborativo de Jogos Digitais}

A metodologia de Desenvolvimento Aberto discutida por \citeonline{thominet2018} pode ser posicionada como um elemento central no ciclo de desenvolvimento de jogos digitais. Segundo o autor, essa estratégia, além de reafirmar um compromisso de transparência com os usuários, incentiva o acesso frequente ao produto, permitindo identificar pontos de melhoria e construir um público mais engajado.

Corroborando essa ideia, \citeonline{vargas2018} analisam como a empresa \textit{Niantic} incorporou, durante atualizações do jogo \textit{Pokémon Go}, conteúdos baseados em demandas identificadas nos comentários deixados pelos consumidores, conseguindo, assim, uma melhor avaliação sobre o desempenho do jogo. Complementarmente a isso, \citeonline{linaker2024} argumentam que a validação contínua com o público é uma excelente estratégia para a redução de riscos e orientação de decisões durante o desenvolvimento.

Complementarmente, \citeonline{alves2019} destacam que o envolvimento dos usuários nas etapas de concepção, prototipagem e teste promove a transição de consumidores para co-autores de tecnologia, ressaltando que o valor da colaboração reside na mediação social e no próprio processo de interação.

% \section{COMUNIDADES ONLINE}

\section{Usabilidade e Experiência do Usuário (UX)}

\textcolor{red}{*Traz os conceitos aqui e traz as heurísticas de Nielsen ou diretrizes de ux e usabilidade que serão adotadas para construir o protótipo (aqui você só menciona as diretrizes e a importância delas, na metodologia vc coloca como seu protótipo atende elas). 
Tente inserir alguma figura ou quadro para ilustrar melhor, ok?}

% \subsection{AVALIAÇÃO DE INTERFACES}

% Sob a perspectiva de Interação Humano-Computador (IHC), a avaliação de sistemas interativos prioriza a qualidade de uso, que analisa o impacto da tecnologia na vida dos usuários e a adequação da interface durante a interação \cite{barbosa2010}. A qualidade de uso articula quatro critérios fundamentais: a usabilidade, que se refere à facilidade de aprendizado e uso da interface, abrangendo fatores como eficácia, eficiência, facilidade de recordação, segurança e satisfação no alcance de objetivos; a experiência do usuário (\textit{User Experience} - UX), que investiga a subjetividade do usuário, englobando as emoções, sentimentos e reações estéticas provocadas durante e após a interação; a acessibilidade, focada na remoção de barreiras na interface para garantir que pessoas com diferentes necessidades ou limitações possam interagir com autonomia; e a comunicabilidade, fundamentada na Engenharia Semiótica, que trata da capacidade da interface de comunicar com clareza a lógica do \textit{design} e as intenções do projetista, permitindo um uso produtivo e criativo da tecnologia \cite{barbosa2010}.

% Segundo \citeonline{pressman2011}, a avaliação do projeto de interface é um processo iterativo no qual um protótipo operacional é submetido a testes com usuários finais e a inspeções de especialistas (como avaliações heurísticas e revisões cognitivas). Essa etapa coleta dados qualitativos, como satisfação e percepção estética, e quantitativos, como tempo de execução e taxa de erros, permitindo refinamentos sucessivos para garantir a usabilidade do sistema.